\textbf{Phần 1. Trắc nghiệm lựa chọn}
\setcounterpageref{ex}{0}

\begin{ex} %[5D7V42]
Dùng máy tính cầm tay để tính giá trị của biểu thức: $48{,}5 \times 12{,}4 - 315{,}8$. Kết quả hiển thị trên màn hình là
\choice
{$275{,}6$}
{\True $285{,}6$}
{$295{,}6$}
{$301{,}4$}
\loigiai{
Thực hiện phép nhân trước: $48{,}5 \times 12{,}4 = 601{,}4$.\\
Sau đó trừ: $601{,}4 - 315{,}8 = 285{,}6$.\\
\textbf{Đáp án của câu hỏi là B}
}
\end{ex}

\begin{ex} %[5D7V42]
Để tìm tỉ số phần trăm của $18$ so với $75$ bằng máy tính cầm tay, ta có thể bấm lần lượt: \fbox{1} \fbox{8} \fbox{$\div$} \fbox{7} \fbox{5} \fbox{\%}. Kết quả màn hình hiển thị là
\choice
{$0{,}24$}
{\True $24$}
{$2{,}4$}
{$240$}
\loigiai{
Ta có: $18 : 75 = 0{,}24 = 24\%$.\\
\textbf{Đáp án của câu hỏi là B}
}
\end{ex}

\begin{ex} %[5D7V42]
Một học sinh cần tính $25 \times 14$ nhưng lỡ tay bấm nhầm thành $25 \times 15$. Để nhận được kết quả đúng mà không phải bấm lại từ đầu, bạn đó chỉ cần bấm trừ đi
\choice
{$1$}
{$10$}
{$14$}
{\True $25$}
\loigiai{
Ta có: $25 \times 15 - 25 \times 14 = 25 \times (15 - 14) = 25$.\\
Do đó chỉ cần bấm trừ đi $25$.\\
\textbf{Đáp án của câu hỏi là D}
}
\end{ex}

\begin{ex} %[5D7V42]
Khi dùng máy tính cầm tay thực hiện phép chia $125 : 8$, màn hình hiển thị kết quả là
\choice
{$15{,}6$}
{\True $15{,}625$}
{$15{,}25$}
{$15{,}75$}
\loigiai{
Ta có: $125 : 8 = 15{,}625$.\\
\textbf{Đáp án của câu hỏi là B}
}
\end{ex}

\begin{ex} %[5D7V42]
Dùng máy tính cầm tay tính giá trị của biểu thức: $(35{,}2 + 14{,}8) \times 2{,}5 : 4$. Kết quả thu được là
\choice
{$30{,}5$}
{\True $31{,}25$}
{$32{,}5$}
{$35$}
\loigiai{
Ta có: $(35{,}2 + 14{,}8) = 50$.\\
$50 \times 2{,}5 = 125$.\\
$125 : 4 = 31{,}25$.\\
\textbf{Đáp án của câu hỏi là B}
}
\end{ex}

\begin{ex} %[5D7V42]
Để tính $18\%$ của $450$ trên máy tính bỏ túi, ta thao tác bấm lần lượt: \fbox{4} \fbox{5} \fbox{0} \fbox{$\times$} \fbox{1} \fbox{8} \fbox{\%}. Kết quả hiển thị là
\choice
{$72$}
{\True $81$}
{$90$}
{$810$}
\loigiai{
Ta có: $450 \times 18\% = 450 \times 0{,}18 = 81$.\\
\textbf{Đáp án của câu hỏi là B}
}
\end{ex}

\begin{ex} %[5D7V42]
Để tính đúng giá trị của biểu thức $100 - 25 \times 3$ trên máy tính bỏ túi thông thường (không có phím ngoặc đơn), ta cần thực hiện thứ tự bấm như thế nào?
\choice
{Bấm lần lượt từ trái sang phải: \fbox{1}\fbox{0}\fbox{0} \fbox{--} \fbox{2}\fbox{5} \fbox{$\times$} \fbox{3} \fbox{=}}
{\True Tính phép nhân trước \fbox{2}\fbox{5} \fbox{$\times$} \fbox{3} \fbox{=}, sau đó bấm \fbox{1}\fbox{0}\fbox{0} \fbox{--} \fbox{7}\fbox{5} \fbox{=}}
{Bấm \fbox{1}\fbox{0}\fbox{0} \fbox{$\times$} \fbox{3} \fbox{--} \fbox{2}\fbox{5} \fbox{=}}
{Bấm phím nào trước cũng được}
\loigiai{
Vì máy tính thông thường thực hiện tính tuần tự từ trái sang phải, nếu bấm liên tiếp sẽ ra $(100 - 25) \times 3 = 225$ là sai quy tắc toán học.\\
Quy tắc đúng là nhân chia trước, cộng trừ sau: $100 - 25 \times 3 = 100 - 75 = 25$.\\
\textbf{Đáp án của câu hỏi là B}
}
\end{ex}

\begin{ex} %[5D7V42]
Dùng máy tính cầm tay tính diện tích hình tròn có bán kính $r = 4{,}5\text{ cm}$ (lấy giá trị $\pi = 3{,}14$). Kết quả là
\choice
{$28{,}26\text{ cm}^2$}
{$56{,}52\text{ cm}^2$}
{\True $63{,}585\text{ cm}^2$}
{$63{,}85\text{ cm}^2$}
\loigiai{
Diện tích hình tròn là:
$3{,}14 \times 4{,}5 \times 4{,}5 = 63{,}585\text{ (cm}^2\text{)}$.\\
\textbf{Đáp án của câu hỏi là C}
}
\end{ex}

\begin{ex} %[5D7V42]
Một món hàng có giá niêm yết $850\,000$ đồng, được giảm giá $12\%$. Người bán hàng bấm trên máy tính: \fbox{8}\fbox{5}\fbox{0}\fbox{0}\fbox{0}\fbox{0} \fbox{--} \fbox{1}\fbox{2} \fbox{\%}. Số tiền khách hàng phải thanh toán là
\choice
{$740\,000\text{ đồng}$}
{\True $748\,000\text{ đồng}$}
{$752\,000\text{ đồng}$}
{$760\,000\text{ đồng}$}
\loigiai{
Số tiền khách phải trả là:
$850\,000 \times (100\% - 12\%) = 850\,000 \times 88\% = 748\,000\text{ (đồng)}$.\\
\textbf{Đáp án của câu hỏi là B}
}
\end{ex}

\begin{ex} %[5D7V42]
Trên máy tính cầm tay, phím chức năng \fbox{M+} thường có tác dụng gì?
\choice
{Xóa bộ nhớ tạm thời của máy}
{\True Cộng giá trị đang hiển thị trên màn hình vào bộ nhớ tạm}
{Trừ giá trị đang hiển thị trên màn hình khỏi bộ nhớ tạm}
{Tắt nguồn máy tính}
\loigiai{
Phím \fbox{M+} dùng để cộng số đang hiển thị vào bộ nhớ tạm thời.\\
\textbf{Đáp án của câu hỏi là B}
}
\end{ex}

\begin{ex} %[5D7V42]
Khi bấm lần lượt các phím: \fbox{2} \fbox{.} \fbox{4} \fbox{$\times$} \fbox{5} \fbox{.} \fbox{5} \fbox{=}, kết quả hiển thị trên màn hình là
\choice
{$12{,}2$}
{\True $13{,}2$}
{$13{,}22$}
{$14{,}2$}
\loigiai{
Ta có: $2{,}4 \times 5{,}5 = 13{,}2$.\\
\textbf{Đáp án của câu hỏi là B}
}
\end{ex}

\begin{ex} %[5D7V42]
Dùng máy tính cầm tay tính trung bình cộng của 5 số: $14{,}2;\ 18{,}5;\ 22{,}1;\ 19{,}4;\ 15{,}8$. Kết quả thu được là
\choice
{$17{,}5$}
{\True $18$}
{$18{,}5$}
{$19$}
\loigiai{
Tổng 5 số là: $14{,}2 + 18{,}5 + 22{,}1 + 19{,}4 + 15{,}8 = 90$.\\
Trung bình cộng là: $90 : 5 = 18$.\\
\textbf{Đáp án của câu hỏi là B}
}
\end{ex}

\textbf{Phần 2. Trắc nghiệm đúng sai}
\setcounterpageref{ex}{0}

\begin{ex} %[5D7V42]
Bạn Hoa dùng máy tính cầm tay để tính giá trị biểu thức $A = 125 \times 0{,}8 + 45 : 1{,}5$.
\choiceTF
{\True Phép tính $125 \times 0{,}8$ cho kết quả bằng $100$}
{\True Phép tính $45 : 1{,}5$ cho kết quả bằng $30$}
{\True Giá trị đúng của biểu thức $A$ bằng $130$}
{Nếu bấm liên tiếp trên máy tính thông thường: $125 \times 0{,}8 + 45 : 1{,}5 =$ thì máy sẽ cho ra đúng kết quả $130$}
\loigiai{
\begin{itemchoice}
\item Đúng. $125 \times 0{,}8 = 100$.
\item Đúng. $45 : 1{,}5 = 30$.
\item Đúng. $A = 100 + 30 = 130$.
\item Sai. Máy tính thông thường thực hiện lần lượt từ trái sang phải: lấy $(100 + 45) : 1{,}5 = 145 : 1{,}5 \approx 96{,}67 
eq 130$.
\end{itemchoice}
}
\end{ex}

\begin{ex} %[5D7V42]
Một người bán hàng dùng máy tính để tính tiền hàng cho khách mua: $3{,}5\text{ kg}$ cam giá $28\,000$ đồng/kg và $2\text{ kg}$ táo giá $45\,000$ đồng/kg.
\choiceTF
{\True Số tiền mua cam là $98\,000$ đồng}
{\True Số tiền mua táo là $90\,000$ đồng}
{\True Tổng số tiền khách phải thanh toán là $188\,000$ đồng}
{Khách đưa tờ $200\,000$ đồng, máy tính bấm $200\,000 - 188\,000 =$ hiển thị số tiền trả lại là $22\,000$ đồng}
\loigiai{
\begin{itemchoice}
\item Đúng. Tiền cam: $3{,}5 \times 28\,000 = 98\,000\text{ (đồng)}$.
\item Đúng. Tiền táo: $2 \times 45\,000 = 90\,000\text{ (đồng)}$.
\item Đúng. Tổng tiền: $98\,000 + 90\,000 = 188\,000\text{ (đồng)}$.
\item Sai. Số tiền trả lại là: $200\,000 - 188\,000 = 12\,000\text{ (đồng)} 
eq 22\,000\text{ đồng}$.
\end{itemchoice}
}
\end{ex}

\textbf{Phần 3. Trắc nghiệm trả lời ngắn}
\setcounterpageref{ex}{0}

\begin{ex} %[5D7V42]
Dùng máy tính cầm tay tính giá trị của biểu thức: $45 \times 18 + 55 \times 18$.
\shortans{1800}
\loigiai{
Ta có: $18 \times (45 + 55) = 18 \times 100 = 1\,800$.
}
\end{ex}

\begin{ex} %[5D7V42]
Dùng máy tính cầm tay tính $25\%$ của số $480$.
\shortans{120}
\loigiai{
Ta có: $480 \times 25\% = 120$.
}
\end{ex}

\begin{ex} %[5D7V42]
Tìm số tự nhiên $x$, biết rằng khi dùng máy tính bấm $x : 15$ thì màn hình cho kết quả là $24$.
\shortans{360}
\loigiai{
Số tự nhiên $x$ là:
$24 \times 15 = 360$.
}
\end{ex}

\begin{ex} %[5D7V42]
Dùng máy tính cầm tay tính giá trị của biểu thức: $(14{,}7 + 25{,}3) \times 15$.
\shortans{600}
\loigiai{
Ta có: $(14{,}7 + 25{,}3) \times 15 = 40 \times 15 = 600$.
}
\end{ex}

\begin{ex} %[5D7V42]
Một đơn hàng có tổng trị giá $400\,000$ đồng, khách hàng được giảm giá $15\%$. Dùng máy tính tính số tiền được giảm theo đơn vị nghìn đồng.
\shortans{60}
\loigiai{
Số tiền được giảm là:
$400\,000 \times 15\% = 60\,000\text{ (đồng)}$.\\
Đổi: $60\,000\text{ đồng} = 60\text{ nghìn đồng}$.
}
\end{ex}

\begin{ex} %[5D7V42]
Dùng máy tính tính giá trị của biểu thức: $3\,600 : (12 \times 3)$.
\shortans{100}
\loigiai{
Ta có: $3\,600 : (12 \times 3) = 3\,600 : 36 = 100$.
}
\end{ex}
